\section{第五日第二讲}

\textbf{免责声明}：我尽力忠实复现这场报告，但对可能引入的任何不确切或遗漏表示歉意。特别是，我发现自己很难复现那些只是口头说出、并未写下的内容。

\textbf{记号说明}：为了把不同记录人的笔记合成一份统一的讲义，我将用 $(-)^{\mathrm{sh}}$ 表示 BZSV 中 shearing（「斜框」）的记号。注意这个记号的含义在报告中并未精确定义，只是大致说过：它把非零的上同调次数转成 $\mathbb G_m$-分次。

\begin{conj}
\footnote{下面这个猜想（粗略版本）在报告中并未完整写出，其中的 $L$ 是我按自己的理解补上的。}
在几何 Langlands 对应下
\[
  \begin{array}{rcl}
    \Shv(\Bun_G)&\xrightarrow{\ (-)_{\mathrm{spec}}\ }&\Shv_N(\Bun_G)\simeq\QC^!(\Loc_{\check G}),\\[2pt]
    P&\mapsto&L.
  \end{array}
\]
对涉及 $\Hom$ 的问题，
\[
  \Hom(P_{\mathrm{spec}},\text{Hecke 特征层})=\Hom(P,\text{「自守形式」})
\]
对上面的谱投影 $(-)_{\mathrm{spec}}$ 不敏感。
\end{conj}

此外还有
\[
  \begin{array}{rcl}
    \bigotimes_{s\in S}\Shv(X_F/G_O)&\xrightarrow{\ \Theta\ }&\Shv(\Bun_G(\Sigma)),\\[2pt]
    \delta_0:=k_{X_O/G_O}\ (\text{基本对象})&\mapsto&P_X=\Theta(\delta_0)\ (\text{周期层}).
  \end{array}
\]
上面的映射与两侧 $H$ 的作用都相容。

在函数层面：
\[
  \begin{array}{rcl}
    C_c^\infty(X_F/G_O)&\xrightarrow{\ \Theta\ }&\Fun(\{\text{在 }s\text{ 处带平凡化的 }G\text{-丛}\})\ni P=\Theta(\delta_0).
  \end{array}
\]

上述内容有一个全分歧版本（下面两行中标出的 $G_F$ 在两侧同时作用）：
\[
  \begin{array}{rcl}
    \Shv(X_F)&\to&\Shv(\Bun_{G,s}),\qquad G_F\text{ 在两侧作用},\\[2pt]
    C_c^\infty(X_F)&\xrightarrow{\ \Theta\ }&\Fun(\Bun_G(\mathbb F_q)),\qquad G_F\text{ 在两侧作用}.
  \end{array}
\]

在极化情形 $\check M=T^*\check X$，我们如下定义相应的 L-层：
\[
  \begin{array}{rcl}
    \QC_{\mathrm{sh}}(\check M/\check G)&\to&\QC^!(\Loc_{\check G}\Sigma),\\[2pt]
    \mathcal O&\mapsto&L_{\check X}.
  \end{array}
\]
注意左端等价于 $\QC^!(\Loc^{\check X}_{\check G}(\text{去心圆盘}))$，而上面 $\mathcal O$ 对应的对象是对偶化层 $\omega$。

\begin{exmp}
$G\circlearrowleft G/H=X$，$X/G=\bullet/H$，
\[
  \Bun_H=\Bun^X_G\xrightarrow{\ \pi\ }\Bun_G,\qquad P_{G/H}=\pi_!k_{\Bun_H},
\]
\[
  \Hom(P_{G/H},F)=\Gamma(\Bun_H,\pi_!F),\qquad P_H(f)=\int_{[H]}f\big|_{[H]},
\]
其中 $f$ 是 $[G]$ 上的自守形式。
\end{exmp}

\begin{exmp}
$G\circlearrowleft G/N=X\circlearrowleft T$，$X/(G\times T)=\bullet/B$，
\[
  \Bun^X_{G\times T}=\Bun_B\ \longrightarrow\ \Bun_G,\qquad \Bun_T.
\]
由这个对应得到 Eisenstein 级数函子
\[
  \mathrm{Eis}:\ \Shv(\Bun_T)\rightleftarrows\Shv(\Bun_G):\ \text{常项（Constant Term）}.
\]
它是由 $P_{G\circlearrowleft G/N\circlearrowleft T}$ 表示的积分变换。这里的对偶是
\[
  G\circlearrowleft G/N\circlearrowleft T\quad\Longleftrightarrow\quad\check G\circlearrowleft\check G/\check N\circlearrowleft\check T
\]
以及
\[
  \begin{array}{ccc}
    \Shv(\Bun_G)&\longleftrightarrow&\QC^!(\Loc_{\check G})\\
    \updownarrow&&\updownarrow\\
    \Shv(\Bun_T)&\longleftrightarrow&\QC^!(\Loc_{\check T})\\[2pt]
    \Bun_B&&\Loc_{\check B}.
  \end{array}
\]
\end{exmp}

\begin{rem}
周期层与 Hecke 特征层相距甚远。在谱侧，这大致就是摩天楼层与 L-层之间的差别。
\end{rem}

\begin{exmp}[群情形]
球簇
\[
  G\times G\circlearrowleft G\quad\Longleftrightarrow\quad\check G\times\check G\circlearrowleft\check G\quad(\text{模去 Chevalley 扭})
\]
对应下列边界条件：
\[
  G\ \big|^{\mathrm{id}}G\ \longrightarrow\ \check G\ \big|^{\mathrm{id}}\check G
\]
以及
\[
  \begin{array}{c}
    \Bun^{X=G}_{G\times G}\simeq\Bun_G\\[1pt]
    {\scriptstyle\Delta}\big\downarrow\\[1pt]
    \Bun_G\times\Bun_G=\Bun_{G\times G}
  \end{array}
\]
其中 $\Delta$ 是对角嵌入，且
\[
  P_{\mathrm{group}}=\Delta_!k,\qquad L_{\mathrm{group}}=\Delta_!\omega.
\]
这编码了 Gaitsgory 的 Miraculous Duality，它说 $\Shv(\Bun_G)$ 有一个版本的 Verdier 对偶，与谱侧的 Serre 对偶相匹配。
\end{exmp}

\begin{exmp}
\[
  \mathbb G_m\circlearrowleft\bullet\ (\text{Neumann})\quad\Longleftrightarrow\quad\mathbb G_m\circlearrowleft\mathbb G_m\ (\text{Dirichlet}).
\]
右端 $\mathbb G_m$ 在 $\mathbb G_m$ 上的作用是（左）乘法\footnote{补充：当然，由于 $\mathbb G_m$ 交换，左乘右乘没有区别。}，而这是一个非球簇的例子。
\end{exmp}

\begin{exmp}
$M=T^*G/\!/_f\check N\circlearrowleft G$（$f\in\mathfrak n^*$），$G\circlearrowleft\bullet$，
\[
  P_\bullet=k_{\Bun_G}\quad\longleftrightarrow\quad L_{\mathrm{Whit}}\ (\text{谱 Whittaker 层}).
\]
Whittaker 归一化：
\[
  P_{\mathrm{Whit}}\quad\longleftrightarrow\quad L_\bullet=\mathcal O_{\Loc_{\check G}}=\omega_{\Loc_{\check G}},
\]
\[
  \Hom(L_\bullet,-)=R\Gamma(\Loc_{\check G},-),\qquad \Hom(P_{\mathrm{Whit}},-)=R\Gamma(\Bun_N,\text{对某个特征标}).
\]
\end{exmp}

\begin{exmp}\footnote{补充：注意下面 $\underline{\mathrm{Gr}}_G$ 带下划线，指的不是仿射 Grassmann 流形 $\mathrm{Gr}_G$，而是 Hecke 栈 $G_O\backslash G_F/G_O$。}
$\Sigma=$「raviolo」$=D\sqcup_{D^*}\!D$。此时
\[
  P_\bullet=k_{\underline{\mathrm{Gr}}_G},\qquad \Hom(P_\bullet,-)=H^*(\underline{\mathrm{Gr}}_G,-)=H^*_{G(O)}(\mathrm{Gr}_G,-),
\]
以及
\[
  c=\mathfrak g/\!/\check G=\mathfrak h/\!/W\simeq\check{\mathfrak h}^*/\!/W,
\]
\[
  \Shv(\Bun_G(\text{raviolo}))\ \leftrightarrow\ \QC_{\mathrm{sh}}(\check{\mathfrak g}^*/\check G),\qquad
  H^*_G(\bullet)\text{-}\mathrm{Mod}=\mathbb C[\check G]\text{-}\mathrm{Mod}\ \simeq\ K^*.
\]
（回到 Whittaker 情形）有
\[
  \check G\backslash T^*\check G/\!/_f\check N\xrightarrow{\ \mu\ }\check G\backslash\check{\mathfrak g}^*\simeq\check G\backslash(\check G\times c).
\]
$\mathfrak{sl}_2^*$ 的情形的图示：Kostant 截面 $\kappa$ 是特征多项式映射 $\check{\mathfrak g}^*/\check G\to c$ 的截面，由 $f+Z\mathfrak g(e)$ 给出。Hitchin 截面是 Kostant 切片的整体类比，它与整体幂零锥横截相交。Whittaker 层就是 Hitchin 截面上的常值层。
\[
  T^*\Bun_G=\Maps\Big(\Sigma,\frac{\mathfrak g^*}{G}\Big/\!_{G_m}\times\Omega^1\Big)\to\Maps\Big(\Sigma,\frac{c}{G_m}\times\Omega^1\Big).
\]
Gaiotto 拉格朗日（子流形）是一个自然态射的像：它由矩映射（模去 $G$）$\mu:M/G\to\mathfrak g^*/G$ 诱导，
\[
  \Bun^M_G=\Maps(\Sigma,M/G)\to\Maps(\Sigma,\mathfrak g^*/G).
\]
\end{exmp}
